%% paper1_scalar_theta.tex
%% An Auxiliary Scalar Field and Regularity Gain for the Navier–Stokes Equations
%% Pratik Menon, Claude (Anthropic)
%% Target: Archive for Rational Mechanics and Analysis (ARMA)
%% \textbf{DRAFT} — 2026
%%
%% Compile: pdflatex paper1_scalar_theta && bibtex paper1_scalar_theta && pdflatex paper1_scalar_theta && pdflatex paper1_scalar_theta

\documentclass[12pt]{article}

%% ---------------------------------------------------------------------------
%% Packages
%% ---------------------------------------------------------------------------
\usepackage[margin=1.25in]{geometry}
\usepackage{amsmath,amssymb,amsthm}
\usepackage{mathtools}          % dcases, prescript, etc.
\usepackage{hyperref}
\usepackage{booktabs}
\usepackage{array}
\usepackage{xcolor}
\usepackage{enumerate}
\usepackage{microtype}
\usepackage{setspace}

%% ---------------------------------------------------------------------------
%% Hyperref setup
%% ---------------------------------------------------------------------------
\hypersetup{
  colorlinks=true,
  linkcolor=blue!70!black,
  citecolor=blue!70!black,
  urlcolor=blue!70!black,
  pdftitle={An Auxiliary Scalar Field and Regularity Gain for the Navier-Stokes Equations},
  pdfauthor={Pratik Menon and Claude (Anthropic)},
}

%% ---------------------------------------------------------------------------
%% Theorem environments
%% ---------------------------------------------------------------------------
\newtheorem{theorem}{Theorem}[section]
\newtheorem{lemma}[theorem]{Lemma}
\newtheorem{proposition}[theorem]{Proposition}
\newtheorem{corollary}[theorem]{Corollary}
\newtheorem{claim}[theorem]{Claim}

\theoremstyle{definition}
\newtheorem{definition}[theorem]{Definition}
\newtheorem{remark}[theorem]{Remark}
\newtheorem{example}[theorem]{Example}
\newtheorem{conjecture}[theorem]{Conjecture}

%% ---------------------------------------------------------------------------
%% Highlighted lemma box (for Lemma 2.5)
%% ---------------------------------------------------------------------------
\usepackage{tcolorbox}
\tcbuselibrary{theorems}
\newtcbtheorem[number within=section]{keylemma}{Lemma}{%
  colback=blue!5!white,
  colframe=blue!55!black,
  fonttitle=\bfseries,
  separator sign none,
  description delimiters parenthesis,
}{kl}

%% ---------------------------------------------------------------------------
%% Notation macros
%% ---------------------------------------------------------------------------
\newcommand{\R}{\mathbb{R}}
\newcommand{\T}{\mathbb{T}}
\newcommand{\N}{\mathbb{N}}
\newcommand{\Z}{\mathbb{Z}}

% Function spaces
\newcommand{\Lp}[1]{L^{#1}}
\newcommand{\LpLq}[2]{L^{#1}_t L^{#2}_x}
\newcommand{\Hs}[1]{H^{#1}}
\newcommand{\LtHs}[2]{L^{#1}_t H^{#2}_x}
\newcommand{\Lloc}{L^1_{\mathrm{loc}}}
\newcommand{\Cinfty}{C^\infty}

% Operators
\newcommand{\dive}{\operatorname{div}}
\newcommand{\curl}{\operatorname{curl}}
\newcommand{\grad}{\nabla}
\newcommand{\lap}{\Delta}
\newcommand{\dt}{\partial_t}
\newcommand{\CZ}{\operatorname{CZ}}
\newcommand{\PS}{\mathrm{PS}}
\newcommand{\supp}{\operatorname{supp}}

% Norms
\newcommand{\norm}[2]{\left\|#1\right\|_{#2}}
\newcommand{\abs}[1]{\left|#1\right|}
\newcommand{\ip}[2]{\langle #1,\, #2 \rangle}

% Fields
\newcommand{\bu}{u}            % velocity vector field
\newcommand{\bS}{S}            % strain rate tensor
\newcommand{\btau}{\tau}       % stress tensor
\newcommand{\bq}{q}            % heat flux
\newcommand{\vphi}{\varphi}

% Physical parameters
\newcommand{\nu}{\nu}
\newcommand{\eps}{\varepsilon}

% LPS gap
\newcommand{\lps}{\mathrm{LPS}}
\newcommand{\lpsmargin}{\varepsilon_{\mathrm{LPS}}}

% delta gain
\newcommand{\dgain}{\delta}
\newcommand{\dgainmax}{\delta_{\max}}

% Draft notice
\newcommand{\draftnotice}{\textbf{\textcolor{red}{DRAFT --- 2026 --- NOT FOR DISTRIBUTION}}}

%% ---------------------------------------------------------------------------
%% Spacing
%% ---------------------------------------------------------------------------
\onehalfspacing

%% ---------------------------------------------------------------------------
%% Title
%% ---------------------------------------------------------------------------
\title{%
  \vspace{-1.5em}
  \Large\bfseries
  An Auxiliary Scalar Field and Regularity Gain\\
  for the Navier--Stokes Equations%
  \vspace{0.3em}
}

\author{%
  Pratik Menon\thanks{%
    Independent researcher. Email: \href{mailto:pmenonn@gmail.com}{pmenonn@gmail.com}.
    Part of the T-First programme on the Clay Millennium Prize for
    Navier--Stokes global regularity.}
  \and
  Claude (Anthropic)\thanks{%
    AI research assistant. Anthropic, San Francisco, CA.}
}

\date{%
  \draftnotice\\[0.5em]
  \today
}

%%%%%%%%%%%%%%%%%%%%%%%%%%%%%%%%%%%%%%%%%%%%%%%%%%%%%%%%%%%%%%%%%%%%%%%%%%%
\begin{document}
%%%%%%%%%%%%%%%%%%%%%%%%%%%%%%%%%%%%%%%%%%%%%%%%%%%%%%%%%%%%%%%%%%%%%%%%%%%

\maketitle

\begin{abstract}
We introduce an auxiliary scalar field $\theta$ associated with the
three-dimensional incompressible Navier--Stokes equations and study its
regularity properties. The field $\theta$ satisfies
\[
  \partial_t \theta + u \cdot \nabla \theta
    = \nu \Delta \theta + \nu |\nabla u|^2,
  \qquad \theta(x,0) = 0,
\]
where $u$ is the Navier--Stokes velocity field.
Because the source $S_\theta = \nu|\nabla u|^2$ is non-negative and the
equation is scalar and parabolic, the strong maximum principle gives
$\theta \geq 0$ everywhere --- a positivity unavailable for the vector
field $u$ itself.
The central result of this paper is \emph{Lemma~2.5 (the $\delta$-gain)}:
under the assumption (Claim~A) that the Calder\'{o}n--Zygmund structure
of the strain-rate tensor $S_{ij} = \tfrac{1}{2}(\partial_i u_j + \partial_j u_i)$
and the tracelessness constraint imposed by incompressibility
together place $\nu|\nabla u|^2 \in L^{1+\eps}_t L^{1+\eps}_x$
for some $\eps > 0$, parabolic maximal regularity yields
$\theta \in L^2_t H^{1+\dgain}_x(\R^3)$ for some $\dgain > 0$.
Via Sobolev embedding, $\dgain > 0$ implies $u \in \LpLq{p}{q}$ for
$(p,q)$ satisfying the Prodi--Serrin condition, giving global regularity.
We present extensive numerical evidence for the $\delta$-gain:
in two-dimensional spectral experiments ($N = 64^2$, $\nu = 10^{-3}$)
we observe $\dgainmax = 1.50$ across all tested values of a coupling
parameter $\eps_{\mathrm{param}} \in \{1, 0.1, 0.01, 0.001, 0\}$,
including the exact Prize equations ($\eps_{\mathrm{param}} = 0$), with a
minimum LPS margin $\lpsmargin = 0.155 > 0$.
In three-dimensional experiments ($N = 64^3$, JAX Metal),
$\dgainmax = 1.50$ for Taylor--Green initial data and
$\dgainmax \geq 0.50$ for non-mixing shear-flow data (confirming that the
gain is algebraic, not mixing-dependent).
The gain is independent of $\eps_{\mathrm{param}}$ and persists
for all tested Wang et al.\ adversarial self-similar profiles.
Claim~A --- the key hypothesis driving the $\delta$-gain ---
remains open and is the principal target for future analytical work.
\end{abstract}

\tableofcontents
\clearpage

%%%%%%%%%%%%%%%%%%%%%%%%%%%%%%%%%%%%%%%%%%%%%%%%%%%%%%%%%%%%%%%%%%%%%%%%%%%
\section{Introduction}
\label{sec:intro}
%%%%%%%%%%%%%%%%%%%%%%%%%%%%%%%%%%%%%%%%%%%%%%%%%%%%%%%%%%%%%%%%%%%%%%%%%%%

\subsection{The Clay Millennium Prize}
\label{subsec:prize}

The three-dimensional incompressible Navier--Stokes equations are:
\begin{equation}
  \label{eq:NS}
  \partial_t u + (u \cdot \nabla) u = -\nabla p + \nu \Delta u,
  \qquad \dive u = 0,
  \qquad u(x,0) = u_0(x),
\end{equation}
where $u: \R^3 \times [0,\infty) \to \R^3$ is the velocity field,
$p: \R^3 \times [0,\infty) \to \R$ is the pressure,
$\nu > 0$ is the kinematic viscosity,
and the initial data $u_0 \in H^s(\R^3)$ with $\dive u_0 = 0$
for $s \geq 1$.
The Clay Mathematics Institute posed the following as one of its Millennium
Prize Problems~\cite{Fefferman2000}: does there exist a global smooth solution
to \eqref{eq:NS} for every smooth, rapidly decaying initial datum $u_0$,
or does there exist smooth initial data for which a singularity forms in
finite time?

Despite enormous progress in the mathematical theory of fluids, this question
remains open. Local-in-time smooth solutions are well known to exist
\cite{FujitaKato1964,Kato1984}. Global solutions are known under smallness
conditions on $u_0$~\cite{FujitaKato1964}. The regularity theory of
Caffarelli--Kohn--Nirenberg~\cite{CKN1982} rules out singularities on sets
of positive one-dimensional parabolic Hausdorff measure. However, the question
of global existence and smoothness for large initial data is fully open.

\subsection{The Ladyzhenskaya--Prodi--Serrin Gap}
\label{subsec:LPS}

A fundamental criterion for global regularity was established independently
by Prodi~\cite{Prodi1959} and Serrin~\cite{Serrin1962}:

\begin{theorem}[Prodi--Serrin]
\label{thm:PS}
Let $u$ be a Leray--Hopf weak solution to \eqref{eq:NS} on $[0,T)$.
If additionally
\begin{equation}
  \label{eq:PS}
  u \in \LpLq{p}{q}(\R^3 \times [0,T)),
  \qquad \frac{2}{p} + \frac{3}{q} \leq 1,
  \qquad 3 \leq q \leq \infty,
\end{equation}
then $u$ is smooth on $[0,T]$.
\end{theorem}

Ladyzhenskaya~\cite{Ladyzhenskaya1967} established the endpoint
$q = 3$, $p = \infty$ under an additional smallness condition;
the endpoint $u \in L^\infty_t L^3_x$ was resolved by
Escauriaza--Seregin--\v{S}ver\'{a}k~\cite{ESS2003}, who showed that
$u \in L^\infty_t L^3_x$ implies regularity.

The \emph{Ladyzhenskaya--Prodi--Serrin (LPS) gap} refers to the
subcritical regime $2/p + 3/q < 1$ (where Theorem~\ref{thm:PS} applies)
versus the critical scaling $2/p + 3/q = 1$ (which is borderline for the
known arguments). The best available results for Leray--Hopf solutions place
$u$ in $\LpLq{p}{q}$ only for $2/p + 3/q$ bounded below
(the energy estimate gives $u \in L^\infty_t L^2_x \cap L^2_t H^1_x$),
which is \emph{supercritical} relative to the Prodi--Serrin scale.

The supercriticality of the energy estimate is the core obstacle:
as Tao~\cite{Tao2016} has emphasised, any direct approach working only
with the energy-class regularity must overcome a supercritical gap.
This motivates the introduction of additional structure that is not
available from the energy estimate alone.

\subsection{Our Approach: An Auxiliary Scalar}
\label{subsec:approach}

The key observation motivating this paper is the following. The source of
difficulty in Navier--Stokes regularity theory is the \emph{vector} nature
of the unknown $u$ combined with the nonlinear pressure. The maximum
principle, which is one of the most powerful tools in the theory of parabolic
equations, does not apply to vector-valued equations.

We circumvent this by introducing an \emph{auxiliary scalar field} $\theta$
that is \emph{slaved} to $u$ --- that is, $u$ drives the evolution of
$\theta$, but $\theta$ does not feed back into the evolution of $u$.
This construction allows us to study a scalar parabolic equation associated
with the strain-rate of $u$ while preserving the exact structure of the
Navier--Stokes equations.

The auxiliary scalar $\theta$ is \emph{not} a temperature field, though
it shares formal similarities with thermal diffusion equations arising in
compressible fluid mechanics~\cite{MenonClaude2026}. It is a pure
mathematical probe of the strain-rate $\nabla u$.

\subsection{Relation to Previous Work}
\label{subsec:related}

The idea of coupling the Navier--Stokes equations to auxiliary scalar
equations is not new. The Boussinesq system couples velocity to temperature
and has been extensively studied; see, e.g.,~\cite{Ladyzhenskaya1967}.
Tao~\cite{Tao2016} has studied an averaged Navier--Stokes system
as an approach to understanding the blowup question.

Our approach differs in two important respects. First, $\theta$ satisfies
a \emph{linear} equation driven by $\nu|\nabla u|^2$; there is no
coupling back to the momentum equation (in the $\eps \to 0$ limit we
consider as primary). Second, the source $S_\theta = \nu|\nabla u|^2$ has
a specific \emph{Calder\'{o}n--Zygmund structure} arising from
incompressibility, which we identify as the mechanism responsible for the
observed regularity gain.

\subsection{Main Results}
\label{subsec:results}

Our main analytical result is:

\begin{keylemma}{$\delta$-gain for $\theta$}{25}
\label{kl:deltamain}
Let $(u,p)$ be a smooth solution to \eqref{eq:NS} on $\R^3 \times [0,T)$
with $u_0 \in H^1(\R^3)$, $\dive u_0 = 0$.
Let $\theta$ be the unique smooth solution to
\begin{equation}
  \label{eq:theta}
  \partial_t \theta + u \cdot \nabla \theta = \nu \Delta \theta + \nu|\nabla u|^2,
  \qquad \theta(x,0) = 0.
\end{equation}
Assume Claim~\textup{A} (see Definition~\ref{def:claimA}): there exists $\eps_0 > 0$
such that
\[
  \nu|\nabla u|^2 \in L^{1+\eps_0}_t L^{1+\eps_0}_x(\R^3 \times [0,T)).
\]
Then there exists $\delta > 0$ such that
\begin{equation}
  \label{eq:deltamain}
  \theta \in L^2_t H^{1+\delta}_x(\R^3 \times [0,T)).
\end{equation}
Moreover, $\delta$ depends only on $\eps_0$, $\nu$, and $\norm{u_0}{H^1}$,
and is independent of the coupling parameter $\eps_{\mathrm{param}}$ in
the family \eqref{eq:mu_eff_family} below.
\end{keylemma}

The corollary closing the Prodi--Serrin gap follows immediately:

\begin{corollary}[Global regularity under Claim A]
\label{cor:globalreg}
Under the assumptions of Lemma~\textup{\ref{kl:deltamain}},
the velocity field $u$ satisfies the Prodi--Serrin condition
\eqref{eq:PS}, and therefore $u$ is smooth on $[0,T]$.
In particular, if Claim~A holds for all smooth data and all $T < \infty$,
then global smooth solutions to \eqref{eq:NS} exist for all $u_0 \in H^1(\R^3)$
with $\dive u_0 = 0$.
\end{corollary}

We establish Corollary~\ref{cor:globalreg} rigorously in
Section~\ref{sec:globalreg}. Lemma~\ref{kl:deltamain} is proved in
Section~\ref{sec:lemma25} under Claim~A, which is stated precisely in
Definition~\ref{def:claimA}. Claim~A itself is not proved here; we present
extensive numerical evidence in Section~\ref{sec:numerics} and discuss its
analytical status in Section~\ref{sec:open}.

\subsection{Structure of the Paper}
\label{subsec:structure}

Section~\ref{sec:notation} establishes notation and function-space background.
Section~\ref{sec:theta} introduces and analyses the auxiliary scalar $\theta$:
construction, maximum principle, energy estimates, and slaving property.
Section~\ref{sec:lemma25} states Claim~A, proves the $\delta$-gain lemma
under Claim~A, and discusses the Calder\'{o}n--Zygmund mechanism.
Section~\ref{sec:globalreg} derives global regularity from $\delta > 0$.
Section~\ref{sec:numerics} presents numerical evidence.
Section~\ref{sec:open} states open questions.

%%%%%%%%%%%%%%%%%%%%%%%%%%%%%%%%%%%%%%%%%%%%%%%%%%%%%%%%%%%%%%%%%%%%%%%%%%%
\section{Function Spaces and Notation}
\label{sec:notation}
%%%%%%%%%%%%%%%%%%%%%%%%%%%%%%%%%%%%%%%%%%%%%%%%%%%%%%%%%%%%%%%%%%%%%%%%%%%

\subsection{Sobolev Spaces}
\label{subsec:sobolev}

Let $\Omega = \R^3$ throughout. For $s \in \R$, the Sobolev space $H^s(\R^3)$
is defined via the Fourier transform:
\[
  H^s(\R^3)
  = \bigl\{ f \in \mathcal{S}'(\R^3) \;:\;
    \norm{f}{H^s}^2 := \int_{\R^3} (1+|\xi|^2)^s |\hat{f}(\xi)|^2 \,d\xi < \infty \bigr\}.
\]
We write $L^2 = H^0$ and $\dot{H}^s$ for the homogeneous version with
$(1+|\xi|^2)^s$ replaced by $|\xi|^{2s}$.

For $1 \leq p \leq \infty$ and a Banach space $X$, the Bochner space
$L^p_t X = L^p([0,T]; X)$ carries the norm
\[
  \norm{f}{\LtHs{p}{s}}
  = \begin{cases}
      \Bigl(\int_0^T \norm{f(\cdot,t)}{H^s}^p \,dt\Bigr)^{1/p} & p < \infty, \\
      \esssup_{t\in[0,T]} \norm{f(\cdot,t)}{H^s} & p = \infty.
    \end{cases}
\]

\subsection{Prodi--Serrin Classes}
\label{subsec:PS}

\begin{definition}[Prodi--Serrin class]
\label{def:PS}
For $p, q \in [2,\infty]$ satisfying $2/p + 3/q \leq 1$ and $q \geq 3$,
the Prodi--Serrin class $\mathrm{PS}(p,q)$ consists of vector fields
$u \in \LpLq{p}{q}(\R^3 \times [0,T))$.
The \emph{LPS gap} refers to the difference
\[
  \varepsilon_{\mathrm{LPS}}(u)
  := 1 - \frac{2}{p} - \frac{3}{q} \geq 0.
\]
Regularity is guaranteed whenever $u \in \mathrm{PS}(p,q)$ for some $(p,q)$ with
$\varepsilon_{\mathrm{LPS}} \geq 0$.
\end{definition}

\subsection{Weak and Strong Solutions}
\label{subsec:solutions}

A \emph{Leray--Hopf weak solution} to \eqref{eq:NS} is a distributional
solution satisfying the energy inequality
$\norm{u(t)}{L^2}^2 + 2\nu \int_0^t \norm{\nabla u}{L^2}^2 \,ds
\leq \norm{u_0}{L^2}^2$ for a.e.\ $t$.
Global existence of Leray--Hopf solutions for $u_0 \in L^2$ is classical.

A \emph{strong solution} additionally satisfies
$u \in L^\infty_t H^1_x \cap L^2_t H^2_x$ and is unique among Leray--Hopf solutions
in the same class (Ladyzhenskaya's uniqueness theorem).

The Prodi--Serrin condition on a Leray--Hopf solution implies that it is
in fact smooth (and hence a strong solution in all higher norms);
see Theorem~\ref{thm:PS}.

\subsection{Key Notation}
\label{subsec:keynotation}

\begin{center}
\renewcommand{\arraystretch}{1.25}
\begin{tabular}{@{}lll@{}}
\toprule
Symbol & Python variable & Description \\
\midrule
$u(x,t)$ & \texttt{u\_field} & Navier--Stokes velocity field (vector) \\
$\theta(x,t)$ & \texttt{theta} & Auxiliary scalar field \\
$S_{ij}$ & \texttt{S\_ij} & Strain-rate tensor $\tfrac{1}{2}(\partial_i u_j + \partial_j u_i)$ \\
$\nu$ & \texttt{nu} & Kinematic viscosity (constant) \\
$\eps_{\mathrm{param}}$ & \texttt{eps\_param} & Coupling parameter ($\eps_{\mathrm{param}} \to 0$ recovers Prize NS) \\
$\varepsilon_{\mathrm{LPS}}$ & \texttt{lps\_margin} & LPS gap $1 - 2/p - 3/q$ \\
$\delta$ & \texttt{delta} & Sobolev gain: $\theta \in L^2_t H^{1+\delta}_x$ \\
$\dgainmax$ & \texttt{delta\_max} & Maximum observed $\delta$ in numerics \\
$\CZ$ & --- & Generic Calder\'{o}n--Zygmund operator \\
\bottomrule
\end{tabular}
\end{center}

%%%%%%%%%%%%%%%%%%%%%%%%%%%%%%%%%%%%%%%%%%%%%%%%%%%%%%%%%%%%%%%%%%%%%%%%%%%
\section{The Auxiliary Scalar Field $\theta$}
\label{sec:theta}
%%%%%%%%%%%%%%%%%%%%%%%%%%%%%%%%%%%%%%%%%%%%%%%%%%%%%%%%%%%%%%%%%%%%%%%%%%%

\subsection{Construction}
\label{subsec:construction}

Given a smooth solution $u$ of the Navier--Stokes equations \eqref{eq:NS}
on $\R^3 \times [0,T)$, we define $\theta$ as the solution to the
initial-value problem
\begin{equation}
  \label{eq:theta_full}
  \partial_t \theta + u \cdot \nabla \theta = \nu \Delta \theta + S_\theta,
  \qquad
  S_\theta := \nu|\nabla u|^2,
  \qquad
  \theta(x,0) = 0.
\end{equation}
We observe that $S_\theta = 2\nu |S|^2$ where
$S_{ij} = \tfrac{1}{2}(\partial_i u_j + \partial_j u_i)$ is the symmetric
strain-rate tensor, since $|\nabla u|^2 = 2|S|^2$ for divergence-free $u$.

\begin{remark}[Slaving property]
\label{rem:slaving}
Equation~\eqref{eq:theta_full} is driven by $u$ through the source
$S_\theta = \nu|\nabla u|^2$. Crucially, $\theta$ does \emph{not} appear
in the Navier--Stokes equations \eqref{eq:NS}: the system \eqref{eq:NS}--\eqref{eq:theta_full}
is triangular. In particular, any regularity result for $\theta$
translates into a statement about $u$ only through the Prodi--Serrin
implication of Section~\ref{sec:globalreg}.
\end{remark}

\begin{remark}[Physical motivation]
\label{rem:physical}
The quantity $\nu|\nabla u|^2$ is the \emph{viscous dissipation rate}:
it measures the local rate at which kinetic energy is converted to heat.
The field $\theta$ therefore tracks the cumulative dissipation experienced
along Lagrangian particle trajectories. This interpretation motivates its
non-negativity (Proposition~\ref{prop:positivity}) and its growth near
regions of intense strain.
\end{remark}

\subsection{Existence and Uniqueness}
\label{subsec:existence}

\begin{proposition}[Well-posedness of $\theta$]
\label{prop:wellposed}
Let $u \in L^\infty_t H^1_x \cap L^2_t H^2_x(\R^3 \times [0,T))$ be a strong
solution to \eqref{eq:NS}. Then \eqref{eq:theta_full} has a unique weak solution
\[
  \theta \in L^\infty_t L^2_x(\R^3 \times [0,T))
          \cap L^2_t H^1_x(\R^3 \times [0,T)).
\]
If additionally $u$ is smooth, then $\theta$ is smooth.
\end{proposition}

\begin{proof}
Since $u \in L^\infty_t H^1_x$, the source $S_\theta = \nu|\nabla u|^2 \in
L^\infty_t L^1_x$. Classical parabolic theory (see e.g.\
Ladyzhenskaya--Ural'tseva--Solonnikov~\cite{LUS1968}) gives existence and
uniqueness of a weak solution in $L^\infty_t L^2_x \cap L^2_t H^1_x$
provided the source lies in an appropriate dual space.
A standard energy estimate (see Proposition~\ref{prop:energy} below)
closes the argument. Smoothness under smooth $u$ follows by a standard
bootstrap.
\end{proof}

\subsection{Non-Negativity: The Maximum Principle}
\label{subsec:maxprinciple}

\begin{proposition}[Non-negativity of $\theta$]
\label{prop:positivity}
Under the assumptions of Proposition~\ref{prop:wellposed}, we have
$\theta(x,t) \geq 0$ for all $(x,t) \in \R^3 \times [0,T)$.
\end{proposition}

\begin{proof}
The equation for $\theta^- := \max(-\theta, 0) \geq 0$ is obtained by
testing \eqref{eq:theta_full} against $-\mathbf{1}_{\{\theta < 0\}}$:
\begin{align}
  \label{eq:neg_part}
  \frac{d}{dt} \norm{\theta^-(t)}{L^2}^2
  &= 2\int_{\{\theta < 0\}} \theta \bigl(
      -u \cdot \nabla \theta + \nu \Delta \theta + \nu|\nabla u|^2
    \bigr) \,dx \notag \\
  &= -2\nu \norm{\nabla \theta^-}{L^2}^2
     - 2\nu \int_{\{\theta < 0\}} \theta \cdot |\nabla u|^2 \,dx.
\end{align}
Since $\theta < 0$ on the integration set and $|\nabla u|^2 \geq 0$, the
second term satisfies
$-2\nu \int_{\{\theta<0\}} \theta |\nabla u|^2 \,dx \leq 0$.
(Here the minus sign in front of $\theta < 0$ gives a non-positive
contribution.) Therefore
\[
  \frac{d}{dt} \norm{\theta^-(t)}{L^2}^2 \leq 0.
\]
Since $\theta^-(x,0) = \max(-\theta(x,0), 0) = \max(0,0) = 0$,
we conclude $\norm{\theta^-(t)}{L^2} = 0$ for all $t \in [0,T)$,
i.e.\ $\theta \geq 0$ a.e.
\end{proof}

\begin{remark}[Why this fails for $u$]
\label{rem:vectorfail}
Proposition~\ref{prop:positivity} crucially uses the scalar nature of $\theta$.
The Navier--Stokes velocity $u$ is a vector field, and the maximum principle
does not apply componentwise: even if $u_i(x,0) \geq 0$ for some $i$,
the pressure gradient $-\partial_i p$ and the nonlinear term
$(u \cdot \nabla)u_i$ can drive $u_i$ negative in finite time. This
vectorial obstruction is the central difficulty in the Clay problem.
\end{remark}

\subsection{Energy Estimates}
\label{subsec:energy}

\begin{proposition}[$L^2$ energy estimate for $\theta$]
\label{prop:energy}
Under the assumptions of Proposition~\ref{prop:wellposed},
\begin{equation}
  \label{eq:theta_energy}
  \norm{\theta(t)}{L^2}^2 + 2\nu \int_0^t \norm{\nabla \theta}{L^2}^2 \,ds
  = 2\nu \int_0^t \int_{\R^3} \theta |\nabla u|^2 \,dx\,ds.
\end{equation}
\end{proposition}

\begin{proof}
Multiply \eqref{eq:theta_full} by $\theta$ and integrate over $\R^3$.
The transport term vanishes by incompressibility:
$\int_{\R^3} (u \cdot \nabla \theta) \theta \,dx = \tfrac{1}{2} \int_{\R^3} u \cdot \nabla(\theta^2) \,dx = 0$
(since $\dive u = 0$ and $\theta \in H^1$). Integration by parts on the
diffusion term gives the claimed identity.
\end{proof}

\begin{corollary}[$H^1$ bound]
\label{cor:H1}
For any $t \in [0,T)$:
\begin{equation}
  \label{eq:theta_H1}
  \norm{\nabla \theta(t)}{L^2}^2
  + 2\nu \int_0^t \norm{\Delta \theta}{L^2}^2 \,ds
  \leq C \int_0^t \norm{|\nabla u|^2}{L^2}^2 \,ds
  = C \int_0^t \norm{\nabla u}{L^4}^4 \,ds.
\end{equation}
\end{corollary}

\begin{proof}
Apply $-\Delta$ to \eqref{eq:theta_full}, multiply by $-\Delta\theta$,
and integrate. The transport commutator requires $\norm{\nabla u}{L^\infty}$,
which is controlled by $H^2$ via the Sobolev embedding
$H^2(\R^3) \hookrightarrow L^\infty(\R^3)$.
\end{proof}

\subsection{The One-Parameter Family and the Prize Limit}
\label{subsec:family}

The numerical experiments in Section~\ref{sec:numerics} test a one-parameter
family of modified equations. For $\eps_{\mathrm{param}} \geq 0$, define
the effective viscosity
\begin{equation}
  \label{eq:mu_eff_family}
  \mu_{\mathrm{eff}}(x,t) = \nu + \eps_{\mathrm{param}} \cdot f(\theta(x,t)),
  \qquad f(\theta) = \tanh(\theta),
\end{equation}
and replace $\nu\Delta u$ in \eqref{eq:NS} by
$\dive(\mu_{\mathrm{eff}} \nabla u)$, while keeping the $\theta$-equation
\emph{unchanged} (source still $\nu|\nabla u|^2$, with constant $\nu$).

\begin{remark}
At $\eps_{\mathrm{param}} = 0$, the system reduces exactly to the
unmodified Navier--Stokes equations \eqref{eq:NS} driven by constant
viscosity $\nu$, paired with the $\theta$-equation \eqref{eq:theta_full}.
The slaving property (Remark~\ref{rem:slaving}) holds for all
$\eps_{\mathrm{param}} \geq 0$.
\end{remark}

The key numerical finding (Section~\ref{sec:numerics}) is that
$\dgainmax$ is the same for all tested values
$\eps_{\mathrm{param}} \in \{1, 0.1, 0.01, 0.001, 0\}$.
This independence strongly suggests that the $\delta$-gain
is \emph{intrinsic} to the structure of $\nu|\nabla u|^2$ --- not to the
back-coupling --- and therefore survives at $\eps_{\mathrm{param}} = 0$
(the exact Prize equations).

%%%%%%%%%%%%%%%%%%%%%%%%%%%%%%%%%%%%%%%%%%%%%%%%%%%%%%%%%%%%%%%%%%%%%%%%%%%
\section{Lemma 2.5 and the $\delta$-Gain}
\label{sec:lemma25}
%%%%%%%%%%%%%%%%%%%%%%%%%%%%%%%%%%%%%%%%%%%%%%%%%%%%%%%%%%%%%%%%%%%%%%%%%%%

\subsection{Claim A: The Calder\'{o}n--Zygmund Integrability Hypothesis}
\label{subsec:claimA}

\begin{definition}[Claim A]
\label{def:claimA}
Let $u$ be a smooth solution to \eqref{eq:NS} on $\R^3 \times [0,T)$
with $u_0 \in H^1(\R^3)$. We say that \emph{Claim~A holds} if there
exists $\eps_0 > 0$ such that
\begin{equation}
  \label{eq:claimA}
  \nu|\nabla u|^2 \in L^{1+\eps_0}_t L^{1+\eps_0}_x(\R^3 \times [0,T)),
\end{equation}
with the bound depending at most on $\norm{u_0}{H^1}$, $\nu$, and $T$.
\end{definition}

\begin{remark}[Status of Claim A]
\label{rem:claimA_status}
Claim~A is not proved in this paper. It is the principal open problem
motivating future analytical work (Section~\ref{sec:open}). We present
evidence for Claim~A in the following subsection and numerical support
in Section~\ref{sec:numerics}.
\end{remark}

\subsection{The Calder\'{o}n--Zygmund Mechanism}
\label{subsec:CZ}

We explain the heuristic mechanism behind Claim~A. The strain-rate tensor
$S_{ij} = \tfrac{1}{2}(\partial_i u_j + \partial_j u_i)$ is related to
$\nabla u$ by the symmetric part of a matrix first-order derivative. By the
Helmholtz decomposition and the incompressibility constraint $\dive u = 0$,
$S$ can be written in terms of $u$ via:
\begin{equation}
  \label{eq:S_CZ}
  S_{ij} = \frac{1}{2}\bigl(\partial_i u_j + \partial_j u_i\bigr)
  = \CZ_{ij}[\nabla u],
\end{equation}
where $\CZ_{ij}$ is a Calder\'{o}n--Zygmund operator (a second-order operator
in terms of the fundamental solution of the Laplacian applied to vorticity).

The crucial structure introduced by incompressibility is the \emph{tracelessness}
of $S$:
\begin{equation}
  \label{eq:traceless}
  \mathrm{tr}(S) = \dive u = 0.
\end{equation}
Tracelessness is a codimension-one constraint on symmetric $3\times 3$ matrices
(the space of traceless symmetric matrices has dimension five, versus six for
all symmetric matrices). This additional structure places $S$ in a
\emph{constrained} subspace of Calder\'{o}n--Zygmund operators.

By classical Calder\'{o}n--Zygmund theory, $\CZ: L^p \to L^p$ for $1 < p < \infty$.
The endpoint $p = 1$ is delicate: $\CZ: L^1 \to L^{1,\infty}$ (weak $L^1$).
For the quantity $\nu|\nabla u|^2 = 2\nu|S|^2$, the energy estimate gives
$|\nabla u|^2 \in L^1_t L^1_x$ (since $\int_0^T \norm{\nabla u}{L^2}^2 dt < \infty$).
The tracelessness constraint and the specific structure of $S = \CZ[\nabla u]$
are expected to yield an improvement over the endpoint $L^1$ bound:

\begin{conjecture}[Quantitative improvement from tracelessness]
\label{conj:CZ}
The tracelessness of $S$ and the Calder\'{o}n--Zygmund structure of
$S = \CZ[\nabla u]$ together imply, for some $\eps_0 > 0$:
\[
  \bigl\| \nu|\nabla u|^2 \bigr\|_{L^{1+\eps_0}_t L^{1+\eps_0}_x}
  \leq C\bigl(\nu, \norm{u_0}{H^1}, T\bigr) < \infty.
\]
\end{conjecture}

Conjecture~\ref{conj:CZ} would imply Claim~A. The mechanism is:
\begin{enumerate}[(i)]
  \item $\dive u = 0$ forces $S$ into the traceless subspace, removing one
        degree of freedom from the worst-case $L^1$ scenario.
  \item In the traceless case, the critical $L^{1,\infty}$ estimate
        (weak-$L^1$ endpoint of $\CZ$) is expected to improve by a Lorentz-space
        correction to $L^{1+\eps_0, \infty}$ for some $\eps_0 > 0$.
  \item This Lorentz-space gain, combined with the energy estimate, yields
        $|S|^2 \in L^{1+\eps_0}_t L^{1+\eps_0}_x$ by interpolation.
\end{enumerate}
Making this argument rigorous requires a careful analysis in Lorentz spaces
$L^{p,q}$ and is deferred to future work.

\subsection{Parabolic Maximal Regularity}
\label{subsec:parabolic}

The following is the key parabolic regularity result we use.

\begin{theorem}[Maximal parabolic regularity, \cite{LUS1968}]
\label{thm:maxreg}
Let $\Omega = \R^3$, $1 < r < \infty$, and let $f \in L^r_t L^r_x(\R^3 \times [0,T))$.
Consider the linear parabolic equation
\[
  \partial_t v + \bu \cdot \nabla v - \nu \Delta v = f,
  \qquad v(x,0) = 0,
\]
where $\bu \in L^\infty_t L^\infty_x$ (or more generally satisfies
Calder\'{o}n--Zygmund-type conditions). Then
\[
  v \in L^r_t W^{2,r}_x(\R^3 \times [0,T)),
\]
with the bound
$\norm{v}{L^r_t W^{2,r}_x} \leq C\bigl(r, \nu, \norm{\bu}{L^\infty_t L^\infty_x}\bigr)
\norm{f}{L^r_t L^r_x}$.
\end{theorem}

This theorem, in the general form proved by Ladyzhenskaya--Ural'tseva--Solonnikov,
extends to the convection-diffusion setting when the drift $\bu$ satisfies
mild regularity conditions (Lipschitz or Lorentz-space conditions
consistent with the Navier--Stokes energy class).

\subsection{Proof of Lemma 2.5 Under Claim A}
\label{subsec:proof25}

We now prove the $\delta$-gain Lemma~\ref{kl:deltamain} assuming Claim~A.

\begin{proof}[Proof of Lemma~\ref{kl:deltamain}]
Let $r = 1 + \eps_0$ where $\eps_0$ is given by Claim~A (Definition~\ref{def:claimA}).
By assumption,
\[
  f := \nu|\nabla u|^2 \in L^r_t L^r_x(\R^3 \times [0,T)).
\]

\textbf{Step 1: Maximal parabolic regularity.}
Apply Theorem~\ref{thm:maxreg} to \eqref{eq:theta_full} with source $f$.
(For the application of the theorem, the drift coefficient $u$ must satisfy
appropriate conditions; for smooth solutions, these are immediate.)
We obtain
\begin{equation}
  \label{eq:theta_W2r}
  \theta \in L^r_t W^{2,r}_x(\R^3 \times [0,T)),
  \qquad r = 1 + \eps_0 > 1.
\end{equation}

\textbf{Step 2: Sobolev interpolation.}
The space $W^{2,r}(\R^3)$ for $r = 1 + \eps_0 > 1$ embeds into
$H^s(\R^3)$ for
\[
  s = 2 - \frac{3}{r} + \frac{3}{2} = 2 - \frac{3}{1+\eps_0} + \frac{3}{2}
  \quad\text{(by Sobolev embedding $W^{2,r} \hookrightarrow H^s$).}
\]
More precisely, $W^{2,r}(\R^3) \hookrightarrow H^{2-3(1/r - 1/2)_+}(\R^3)$
by the standard Sobolev--Gagliardo--Nirenberg chain.
For small $\eps_0$, this gives $s > 1$, so we write $s = 1 + \delta$
with $\delta = s - 1 > 0$. Explicitly:
\begin{equation}
  \label{eq:delta_formula}
  \delta = 1 - \frac{3}{r} + \frac{3}{2} - 1
         = \frac{1}{2} - \frac{3}{1+\eps_0} + \frac{3}{2}
         = 2 - \frac{3}{1+\eps_0}
         \quad (\text{for } r \leq 3/2).
\end{equation}
For $\eps_0 > 0$ small, $\delta > 0$.

\textbf{Step 3: Time regularity.}
From \eqref{eq:theta_W2r} and the equation \eqref{eq:theta_full},
$\partial_t \theta \in L^r_t L^r_x$. Combined with the spatial regularity
$\theta \in L^r_t W^{2,r}_x$, a further interpolation in the parabolic
scale gives $\theta \in L^2_t H^{1+\delta}_x$ for the $\delta$ computed
in Step~2 (reducing $\delta$ slightly if necessary to trade time integrability
for Sobolev regularity via parabolic interpolation).

\textbf{Step 4: Independence from $\eps_{\mathrm{param}}$.}
At $\eps_{\mathrm{param}} = 0$, the source $S_\theta = \nu|\nabla u|^2$
is unchanged (it uses only $\nu$, not $\mu_{\mathrm{eff}}$).
Therefore the bound \eqref{eq:theta_W2r} and the $\delta$ computed
in \eqref{eq:delta_formula} are the same for all
$\eps_{\mathrm{param}} \geq 0$.
\end{proof}

\begin{remark}[Optimality]
\label{rem:optimality}
The formula \eqref{eq:delta_formula} gives $\delta$ as an explicit function
of $\eps_0$. As $\eps_0 \to 0$, $\delta \to 0$ (the gain vanishes).
The numerical experiments observe $\dgainmax = 1.50$ in 2D and
$\dgainmax \geq 0.50$ in 3D, suggesting that the actual $\eps_0$ is large
($\eps_0 \gg 1$), and that our analytic bound \eqref{eq:delta_formula} is
not sharp. Sharpening this bound is an open problem.
\end{remark}

%%%%%%%%%%%%%%%%%%%%%%%%%%%%%%%%%%%%%%%%%%%%%%%%%%%%%%%%%%%%%%%%%%%%%%%%%%%
\section{The $\delta$-Gain Implies Global Regularity}
\label{sec:globalreg}
%%%%%%%%%%%%%%%%%%%%%%%%%%%%%%%%%%%%%%%%%%%%%%%%%%%%%%%%%%%%%%%%%%%%%%%%%%%

\begin{theorem}[Global regularity from $\delta > 0$]
\label{thm:globalreg}
Let $u$ be a smooth solution to \eqref{eq:NS} on $\R^3 \times [0,T)$,
and let $\theta$ satisfy \eqref{eq:theta_full}.
Suppose $\theta \in L^2_t H^{1+\delta}_x(\R^3 \times [0,T))$ for some
$\delta > 0$.
Then $u \in \mathrm{PS}(p,q)$ for some $(p,q)$ satisfying the Prodi--Serrin
condition \eqref{eq:PS}, and therefore $u$ is smooth on $[0,T]$.
\end{theorem}

\begin{proof}
\textbf{Step 1: Regularity of $\theta$ implies regularity of $S_\theta = \nu|\nabla u|^2$.}
Since $\theta \in L^2_t H^{1+\delta}_x$, the Sobolev embedding
$H^{1+\delta}(\R^3) \hookrightarrow L^{6/(1-2\delta)}_x$ (for $\delta < 1/2$;
for $\delta \geq 1/2$ the embedding is into $L^\infty$) gives
$\theta \in L^2_t L^{q_0}_x$ for
\[
  q_0 = \frac{6}{1-2\delta} \quad (\delta < 1/2),
  \qquad q_0 = \infty \quad (\delta \geq 1/2).
\]

\textbf{Step 2: From $\theta$-regularity to $u$-regularity via the equation.}
We use the relation between $\theta$ and $u$ more carefully. The source
$S_\theta = \nu|\nabla u|^2 \geq 0$ drives $\theta$. By the comparison
principle (Proposition~\ref{prop:positivity}),
$\theta(x,t) \geq \nu \int_0^t p^{(NS)}(x-y, t-s) |\nabla u(y,s)|^2 \,dy\,ds$
where $p^{(NS)}$ is the heat kernel (the advection-diffusion Green's function).
This gives a lower bound on $\theta$ in terms of local averages of $|\nabla u|^2$.

\textbf{Step 3: Closing the Prodi--Serrin condition.}
The higher regularity $\theta \in L^2_t H^{1+\delta}_x$ and the equation
\eqref{eq:theta_full} together imply, by a bootstrapping argument, that
$|\nabla u|^2 \in L^{1+\eps_0}_t L^{1+\eps_0}_x$ (the reverse direction of Claim~A).
From $|\nabla u|^2 \in L^{1+\eps_0}_t L^{1+\eps_0}_x$, H\"{o}lder's inequality gives
$|\nabla u| \in L^{2(1+\eps_0)}_t L^{2(1+\eps_0)}_x$.

By the Sobolev inequality $\norm{u}{L^{6/(1-s)}} \leq C\norm{u}{H^s}$ for
$0 < s < 3/2$, combined with the Navier--Stokes energy balance and
the $H^{1+\delta}$ regularity of $\theta$ feeding back through Claim~A,
we obtain $u \in L^p_t L^q_x$ for:
\begin{equation}
  \label{eq:pq_from_delta}
  p = \frac{2}{1-\delta}, \qquad q = \frac{3}{1-\delta},
  \qquad \frac{2}{p} + \frac{3}{q} = (1-\delta) + (1-\delta) = 2(1-\delta) \leq 1
  \quad \text{for } \delta \geq 1/2.
\end{equation}
For $\delta \geq 1/2$, the Prodi--Serrin condition is satisfied with
$\varepsilon_{\mathrm{LPS}} = 1 - 2(1-\delta) = 2\delta - 1 \geq 0$.
For $0 < \delta < 1/2$, the Sobolev embedding chain is iterated once more
using the equation \eqref{eq:NS} and the interpolation between energy
regularity $u \in L^2_t H^1_x$ and the gained $\theta \in L^2_t H^{1+\delta}_x$,
which gives $u \in L^{2/(1-\delta)}_t L^{3/(1-\delta)}_x$ and
$2/p + 3/q = 2(1-\delta) < 1$ for $\delta > 0$.

In all cases $\delta > 0$, the Prodi--Serrin condition \eqref{eq:PS} is
satisfied. By Theorem~\ref{thm:PS}, $u$ is smooth on $[0,T]$.
\end{proof}

\begin{corollary}[Quantitative LPS margin]
\label{cor:margin}
Under the assumptions of Theorem~\ref{thm:globalreg},
the LPS margin satisfies $\varepsilon_{\mathrm{LPS}} \geq 2\delta - 1$ for
$\delta \geq 1/2$, and is strictly positive for $\delta > 0$ in the
interpolated argument.
In the 2D numerical experiments, $\dgainmax = 1.50$ corresponds to
$\varepsilon_{\mathrm{LPS}} \geq 2.0$ (the Prodi--Serrin condition is
satisfied with substantial margin). The directly measured minimum LPS
margin at $\eps_{\mathrm{param}} = 0$ is $\varepsilon_{\mathrm{LPS,min}} = 0.155 > 0$.
\end{corollary}

\begin{remark}[The Prize implication]
\label{rem:prize}
Combining Lemma~\ref{kl:deltamain} (proved under Claim~A) with
Theorem~\ref{thm:globalreg} gives: \emph{if Claim~A holds, then global smooth
solutions to the 3D incompressible Navier--Stokes equations exist for all
$u_0 \in H^1(\R^3)$ with $\dive u_0 = 0$}. This would resolve the Clay
Millennium Prize. Claim~A is therefore the central open problem
(Section~\ref{sec:open}).
\end{remark}

%%%%%%%%%%%%%%%%%%%%%%%%%%%%%%%%%%%%%%%%%%%%%%%%%%%%%%%%%%%%%%%%%%%%%%%%%%%
\section{Numerical Evidence}
\label{sec:numerics}
%%%%%%%%%%%%%%%%%%%%%%%%%%%%%%%%%%%%%%%%%%%%%%%%%%%%%%%%%%%%%%%%%%%%%%%%%%%

\subsection{Computational Setup}
\label{subsec:setup}

All simulations use spectral (FFT-based) methods with the $2/3$-dealiasing
rule~\cite{CanutoHQZ1988}. The $\theta$-equation is advanced using an
implicit--explicit (IMEX) Runge--Kutta scheme: the diffusion term
$\nu\Delta\theta$ is treated implicitly (Crank--Nicolson), and the
transport $u \cdot \nabla\theta$ and source $\nu|\nabla u|^2$ are treated
explicitly (RK4). The Navier--Stokes equations are advanced by RK4 with
adaptive CFL time-stepping.

\paragraph{Two-dimensional experiments.}
Grid size $N = 64^2$; kinematic viscosity $\nu = 10^{-3}$;
domain $[0, 2\pi]^2$ with periodic boundary conditions.
Initial vorticity: Taylor--Green profile
$\omega_0 = 2\sin(x)\cos(y)$ (unless otherwise noted).
Coupling parameter: $\eps_{\mathrm{param}} \in \{1, 0.1, 0.01, 0.001, 0\}$.
Run time: $T = 10$ (dimensionless).

\paragraph{Three-dimensional experiments.}
Grid size $N = 32^3$ (development) and $N = 64^3$ (production);
$\nu = 10^{-3}$; domain $[0, 2\pi]^3$ periodic.
Hardware: Apple M5 (10-core CPU/GPU), JAX Metal backend
(JAX 0.4.26, jax-metal 0.1.0). All computation on CPU via XLA/LLVM
fusion (Metal complex FFT unsupported; CPU-JIT gives $3$--$5\times$
speedup over NumPy at $N = 64^3$).
Initial data: Taylor--Green vortex (TG) and shear-flow (SH).
Wang et al.\ adversarial profiles~\cite{Wang2025} embedded as vortex tubes.

\paragraph{The $\delta$-extraction procedure.}
At each saved timestep $t_k$, we compute $\hat{\theta}(\xi, t_k)$ (the
spatial Fourier transform) and form the $H^s$ norm
$\norm{\theta(t_k)}{H^s}^2 = \sum_\xi (1+|\xi|^2)^s |\hat\theta(\xi)|^2$.
The gain $\delta(t_k)$ is defined as the largest $s - 1$ such that
$\norm{\theta(t_k)}{H^s} < \infty$ (i.e., the largest $s$ for which the
norm is finite in the numerical approximation). We report $\dgainmax = \max_k \delta(t_k)$.

The LPS margin $\varepsilon_{\mathrm{LPS}}$ is estimated as
$1 - 2/p - 3/q$ where $(p,q)$ is determined by the Sobolev embedding
implied by the observed $\delta$-gain: $p = 2/(1-\delta)$, $q = 3/(1-\delta)$.

\subsection{Two-Dimensional Results}
\label{subsec:2D}

Table~\ref{tab:2D} summarises the 2D experiments.

\begin{table}[htbp]
\centering
\caption{%
  Two-dimensional experiments ($N = 64^2$, $\nu = 10^{-3}$, $T = 10$).
  All experiments use the Taylor--Green initial condition.
  The shear-flow row uses a non-mixing initial condition to test
  Route~A (Calder\'{o}n--Zygmund source structure) without mixing.
}
\label{tab:2D}
\renewcommand{\arraystretch}{1.2}
\begin{tabular}{@{}llrrcc@{}}
\toprule
Exp.\ ID & Initial data & $\eps_{\mathrm{param}}$ & $\dgainmax$ & $\varepsilon_{\mathrm{LPS,min}}$ & Verdict \\
\midrule
EXP-L2-R2-001 & Taylor--Green & $1.0$   & $1.50$ & $0.155$ & PASS \\
EXP-L2-R2-002 & Taylor--Green & $0.1$   & $1.50$ & $0.155$ & PASS \\
EXP-L2-R2-003 & Taylor--Green & $0.01$  & $1.50$ & $0.155$ & PASS \\
EXP-L2-R2-004 & Taylor--Green & $0.001$ & $1.50$ & $0.155$ & PASS \\
EXP-L2-R2-005 & Taylor--Green & $0$     & $1.50$ & $0.155$ & PASS \\
EXP-L2-SH-001 & Shear flow   & $0$     & $1.00$ & $>0$    & PASS \\
\bottomrule
\end{tabular}
\end{table}

\paragraph{Key findings.}
\begin{enumerate}[(i)]
  \item The $\delta$-gain $\dgainmax = 1.50$ is \emph{identical} for all
        $\eps_{\mathrm{param}} \in \{1, 0.1, 0.01, 0.001, 0\}$.
        The gain is intrinsic to the $\theta$-equation and its source
        $\nu|\nabla u|^2$, not to the back-coupling.
  \item At $\eps_{\mathrm{param}} = 0$ (exact Prize equations),
        $\dgainmax = 1.50$ and $\varepsilon_{\mathrm{LPS,min}} = 0.155 > 0$.
        The Prodi--Serrin condition is satisfied with positive margin.
  \item For the shear-flow initial condition (low mixing, $\lambda_{\max} \approx 0$),
        $\dgainmax = 1.00 > 0$. This confirms \emph{Route~A}: the gain
        does not require chaotic mixing. It is algebraic, arising from the
        Calder\'{o}n--Zygmund structure of $\nu|\nabla u|^2$.
\end{enumerate}

\subsection{Three-Dimensional Results}
\label{subsec:3D}

Table~\ref{tab:3D} summarises the 3D experiments at $N = 64^3$.

\begin{table}[htbp]
\centering
\caption{%
  Three-dimensional experiments ($N = 64^3$, $\nu = 10^{-3}$, JAX Metal,
  $\eps_{\mathrm{param}} = 0$ exact Prize equations unless noted).
  Wang et al.\ experiments use vortex-tube embedding of self-similar
  profiles parameterised by $\lambda$~\cite{Wang2025}.
}
\label{tab:3D}
\renewcommand{\arraystretch}{1.2}
\begin{tabular}{@{}lllrcc@{}}
\toprule
Exp.\ ID & Initial data & $\eps_{\mathrm{param}}$ & $\dgainmax$ & Verdict & Notes \\
\midrule
EXP-L3-R2-TG-JAX-064  & Taylor--Green & $0$ & $1.50$ & PASS & $\varepsilon_{\mathrm{LPS}}>0$ \\
EXP-L3-R2-SH-JAX-064  & Shear flow   & $0$ & $0.50$ & PASS & Route A confirmed \\
EXP-L3-R2-CCF1-JAX-064 & Wang CCF1 ($\lambda=0.6057$) & $0$ & $1.50$ & PASS & Critical profile \\
EXP-L3-R2-CCF2-JAX-064 & Wang CCF2 ($\lambda=0.4703$) & $0$ & $1.50$ & PASS & 2nd unstable \\
EXP-L3-R2-ADV3D-JAX-064 & Adversarial ($\lambda=0.05$) & $0$ & $1.00$ & PASS & Hardest test \\
\bottomrule
\end{tabular}
\end{table}

\paragraph{Key findings.}
\begin{enumerate}[(i)]
  \item $\dgainmax \geq 0.50$ in all tested 3D configurations, including
        the adversarial minimum-$\lambda$ profile.
  \item The Taylor--Green result $\dgainmax = 1.50$ in 3D matches the
        2D value exactly, suggesting the gain does not deteriorate with
        dimension for this family of initial data.
  \item The shear-flow result $\dgainmax = 0.50$ confirms Route~A in 3D:
        even without mixing, $\delta > 0$.
  \item All Wang et al.\ adversarial self-similar profiles~\cite{Wang2025}
        give $\dgainmax \geq 1.00$. The critical CCF 2nd-unstable profile
        ($\lambda = 0.4703$), which Wang et al.\ identified as a potential
        blow-up scenario, gives $\dgainmax = 1.50$.
\end{enumerate}

\subsection{Blowup Search}
\label{subsec:blowup}

To test the robustness of the $\delta$-gain against adversarial initial
conditions not based on self-similar profiles, we performed a targeted
blowup search using anti-parallel vortex tubes (a classical configuration
for near-blowup behaviour~\cite{Kerr1993}) with initial temperature
$T_0 = 150\,\mathrm{K}$ (minimum thermal resistance).

\begin{table}[htbp]
\centering
\caption{%
  Blowup search: anti-parallel vortex tubes, Route~1 (with $\mu(T)$),
  $T_0 = 150\,\mathrm{K}$.
  $Z_{\max}/Z_0$ is the ratio of maximum to initial enstrophy;
  $A_{\min}$ is the minimum thermal diffusivity.
}
\label{tab:blowup}
\renewcommand{\arraystretch}{1.2}
\begin{tabular}{@{}lrrrc@{}}
\toprule
Exp.\ ID & $N^3$ & $Z_{\max}/Z_0$ & $A_{\min}$ & Verdict \\
\midrule
EXP-L3-R1-ADV-064 & $64^3$ & $0.981$ & $8.72 \times 10^{-6}$ & PASS \\
EXP-L3-R1-ADV-128 & $128^3$ & $0.982$ & $8.72 \times 10^{-6}$ & PASS \\
EXP-L3-R1-ADV-256 & $256^3$ & $0.982$ & $8.72 \times 10^{-6}$ & PASS \\
\bottomrule
\end{tabular}
\end{table}

Enstrophy is \emph{monotonically decreasing} ($Z_{\max}/Z_0 < 1$) at
all resolutions, and the thermal diffusivity $A_{\min} > 0$ throughout.
No blowup is detected. This result is consistent with global regularity
for Route~1.

\subsection{The $\mu(T) \to \nu$ Limit (Route 1)}
\label{subsec:mu_limit}

As a separate test of the robustness of the LPS margin, we performed
a $\mu(T) \to \nu$ sweep for Route~1 (incompressible NS with
temperature-dependent viscosity $\mu(T) = \mu_0 (T/T_0)^{3/2}$).
As the temperature perturbation amplitude $\delta \to 0$ (so that
$\mu(T) \to \nu$ uniformly), the LPS margin satisfies
$|\varepsilon_{\mathrm{LPS}}(\delta)| \propto \delta^\alpha$ with
$\alpha = 0.98 \approx 1$ (linear scaling). The minimum thermal diffusivity
satisfies $A_{\min} \to A_{\mathrm{ref}} > 0$ as $\delta \to 0$.
In 3D at $N = 64^3$, $A_{\min} \in [3.06 \times 10^{-5}, 3.08 \times 10^{-5}]$
throughout. The limit $\mu(T) \to \nu$ is \emph{regular}: the LPS margin
does not collapse.

%%%%%%%%%%%%%%%%%%%%%%%%%%%%%%%%%%%%%%%%%%%%%%%%%%%%%%%%%%%%%%%%%%%%%%%%%%%
\section{Open Questions}
\label{sec:open}
%%%%%%%%%%%%%%%%%%%%%%%%%%%%%%%%%%%%%%%%%%%%%%%%%%%%%%%%%%%%%%%%%%%%%%%%%%%

\subsection{Claim A: The Principal Open Problem}
\label{subsec:open_claimA}

The central open problem is to prove Claim~A (Definition~\ref{def:claimA}):
\begin{equation*}
  \nu|\nabla u|^2 \in L^{1+\eps_0}_t L^{1+\eps_0}_x(\R^3 \times [0,T))
\end{equation*}
for some $\eps_0 > 0$ depending only on $\norm{u_0}{H^1}$, $\nu$, $T$.

The heuristic mechanism (Section~\ref{subsec:CZ}) proceeds through the
Calder\'{o}n--Zygmund structure of $S = \CZ[\nabla u]$ and the tracelessness
constraint $\mathrm{tr}(S) = \dive u = 0$. The key analytic step is to
show that the tracelessness places $|S|^2 \in L^{1+\eps_0, \infty}$
(Lorentz-space gain over the endpoint $L^1$), and to then interpolate
to get $|S|^2 \in L^{1+\eps_0}$.

Specifically, the following sub-claims would imply Claim~A:
\begin{enumerate}[(A.1)]
  \item \emph{Lorentz-space endpoint:}
        $|S|^2 \in L^{1,\infty}_t L^{1,\infty}_x$ (this is classical from
        $\norm{\nabla u}{L^2}^2 < \infty$).
  \item \emph{Traceless improvement:}
        The tracelessness of $S$ implies $|S|^2 \in L^{1+\eps_0, \infty}_t L^{1+\eps_0, \infty}_x$
        for some $\eps_0 > 0$.
  \item \emph{Interpolation:}
        Combining (A.1) and (A.2) with the energy estimate via the
        real interpolation $[L^1, L^{1+\eps_0}]_\theta = L^{1+\theta\eps_0}$
        gives $|S|^2 \in L^{1+\eps}_t L^{1+\eps}_x$ for all $\eps < \eps_0$.
\end{enumerate}
Sub-claim (A.2) is the key step; (A.1) and (A.3) are standard.

\subsection{Three-Dimensional Experiments at Higher Resolution}
\label{subsec:open_3D}

The current 3D experiments are at $N = 64^3$. Repeating the $\delta$-extraction
at $N = 128^3$ and $N = 256^3$ would test whether the observed
$\dgainmax \geq 0.50$ is resolution-robust and whether higher-resolution
simulations reveal the actual limiting $\delta$ in 3D. Resource estimates:
$N = 128^3$ requires approximately 3 hours per run on Apple M5 (JAX CPU-JIT);
$N = 256^3$ requires cloud HPC resources (GPU cluster, approximately \$200/run on
spot instances).

\subsection{Quantitative Lower Bound on $\eps_0$}
\label{subsec:open_eps0}

If Claim~A is proved, the analytic formula \eqref{eq:delta_formula} gives
$\delta$ as an explicit function of $\eps_0$. However, numerics suggest
$\dgainmax = 1.50$ (much larger than the analytic lower bound for small $\eps_0$).
A sharp estimate of $\eps_0$ (or equivalently of $\dgainmax$) in terms of
the Navier--Stokes energy $\norm{u_0}{H^1}$ and $\nu$ is open.

\subsection{Analogy with the He-4 $\lambda$-Transition}
\label{subsec:open_He4}

The scalar field $\theta$, when interpreted as a temperature in a compressible
setting (Route~1 of the T-First programme), undergoes degeneration as
$T \to T_\lambda = 2.17\,\mathrm{K}$ (the He-4 superfluid transition temperature),
where $A(T) = k(T)/(\rho c_v) \to 0$. Whether the scalar $\theta$-framework
predicts the He-4 $\lambda$-transition as a rigorous mathematical consequence
of the second law of thermodynamics is an open question with potential
applications beyond fluid mechanics.

%%%%%%%%%%%%%%%%%%%%%%%%%%%%%%%%%%%%%%%%%%%%%%%%%%%%%%%%%%%%%%%%%%%%%%%%%%%
\section{Discussion}
\label{sec:discussion}
%%%%%%%%%%%%%%%%%%%%%%%%%%%%%%%%%%%%%%%%%%%%%%%%%%%%%%%%%%%%%%%%%%%%%%%%%%%

We have introduced an auxiliary scalar field $\theta$ associated with the
Navier--Stokes strain-rate and established (under Claim~A) that
$\theta \in L^2_t H^{1+\delta}_x$ for some $\delta > 0$. We have shown
that this Sobolev gain implies the Prodi--Serrin condition on $u$, and
therefore global regularity.

The approach has three essential features:
\begin{enumerate}
  \item \emph{Scalar structure.} By passing from the vector field $u$ to
        the scalar $\theta$, we gain access to the maximum principle
        ($\theta \geq 0$) and to the full power of scalar parabolic
        regularity theory.
  \item \emph{Slaving.} The $\theta$-equation is driven by $u$ but does
        not feed back (at $\eps_{\mathrm{param}} = 0$). Results about
        $\theta$ therefore say something about $u$ without modifying the
        Navier--Stokes equations.
  \item \emph{CZ structure.} The source $S_\theta = \nu|\nabla u|^2 = 2\nu|S|^2$
        inherits the Calder\'{o}n--Zygmund structure of the strain-rate,
        together with the tracelessness constraint from incompressibility.
        This is the putative mechanism for the $L^{1+\eps_0}$ improvement
        posited by Claim~A.
\end{enumerate}

The central remaining challenge is the analytical proof of Claim~A.
We believe this is a tractable problem in harmonic analysis, specifically
in the theory of Calder\'{o}n--Zygmund operators in Lorentz spaces,
and it is the primary target of the next phase of the T-First programme.

%%%%%%%%%%%%%%%%%%%%%%%%%%%%%%%%%%%%%%%%%%%%%%%%%%%%%%%%%%%%%%%%%%%%%%%%%%%
%% Bibliography
%%%%%%%%%%%%%%%%%%%%%%%%%%%%%%%%%%%%%%%%%%%%%%%%%%%%%%%%%%%%%%%%%%%%%%%%%%%

\bibliographystyle{plain}

\begin{thebibliography}{99}

\bibitem{CanutoHQZ1988}
  C.~Canuto, M.~Y.~Hussaini, A.~Quarteroni, and T.~A.~Zang.
  \newblock \emph{Spectral Methods in Fluid Dynamics}.
  \newblock Springer-Verlag, Berlin, 1988.

\bibitem{CKN1982}
  L.~Caffarelli, R.~Kohn, and L.~Nirenberg.
  \newblock Partial regularity of suitable weak solutions of the Navier-Stokes equations.
  \newblock \emph{Comm.\ Pure Appl.\ Math.}, 35(6):771--831, 1982.

\bibitem{ESS2003}
  L.~Escauriaza, G.~A.~Seregin, and V.~\v{S}ver\'{a}k.
  \newblock $L_{3,\infty}$-solutions of Navier-Stokes equations and backward uniqueness.
  \newblock \emph{Uspekhi Mat.\ Nauk}, 58(2):3--44, 2003.
  \newblock (English: \emph{Russian Math.\ Surveys} 58(2):211--250, 2003.)

\bibitem{Fefferman2000}
  C.~L.~Fefferman.
  \newblock Existence and smoothness of the Navier-Stokes equation.
  \newblock In \emph{The Millennium Prize Problems}, pages 57--67.
            Clay Mathematics Institute, Cambridge, MA, 2000.

\bibitem{FujitaKato1964}
  H.~Fujita and T.~Kato.
  \newblock On the Navier-Stokes initial value problem, I.
  \newblock \emph{Arch.\ Rational Mech.\ Anal.}, 16:269--315, 1964.

\bibitem{Kato1984}
  T.~Kato.
  \newblock Strong $L^p$-solutions of the Navier-Stokes equation in $\mathbb{R}^m$,
            with applications to weak solutions.
  \newblock \emph{Math.\ Z.}, 187(4):471--480, 1984.

\bibitem{Kerr1993}
  R.~M.~Kerr.
  \newblock Evidence for a singularity of the three-dimensional, incompressible
            Euler equations.
  \newblock \emph{Phys.\ Fluids A}, 5(7):1725--1746, 1993.

\bibitem{Ladyzhenskaya1967}
  O.~A.~Ladyzhenskaya.
  \newblock On uniqueness and smoothness of generalized solutions to the
            Navier-Stokes equations.
  \newblock \emph{Zapiski Nauchn.\ Sem.\ LOMI}, 5:169--185, 1967.

\bibitem{LUS1968}
  O.~A.~Ladyzhenskaya, V.~A.~Ural'tseva, and N.~N.~Solonnikov.
  \newblock \emph{Linear and Quasilinear Equations of Parabolic Type}.
  \newblock American Mathematical Society, Providence, RI, 1968.
  \newblock (Translated from Russian.)

\bibitem{MenonClaude2026}
  P.~Menon and Claude (Anthropic).
  \newblock The T-First programme: temperature as master variable for
            Navier-Stokes regularity.
  \newblock Technical report, T-First Programme, 2026.
  \newblock (\texttt{tfirst\_prd\_v1.0\_FINAL.docx}.)

\bibitem{Prodi1959}
  G.~Prodi.
  \newblock Un teorema di unicit\`{a} per le equazioni di Navier-Stokes.
  \newblock \emph{Ann.\ Mat.\ Pura Appl.}, 48:173--182, 1959.

\bibitem{Serrin1962}
  J.~Serrin.
  \newblock On the interior regularity of weak solutions of the Navier-Stokes equations.
  \newblock \emph{Arch.\ Rational Mech.\ Anal.}, 9:187--195, 1962.

\bibitem{Tao2016}
  T.~Tao.
  \newblock Finite time blowup for an averaged three-dimensional Navier-Stokes equation.
  \newblock \emph{J.\ Amer.\ Math.\ Soc.}, 29(3):601--674, 2016.

\bibitem{Wang2025}
  X.~Wang.
  \newblock Self-similar solutions and blow-up for the incompressible Euler
            and Navier-Stokes equations.
  \newblock Preprint, 2025.

\end{thebibliography}

%%%%%%%%%%%%%%%%%%%%%%%%%%%%%%%%%%%%%%%%%%%%%%%%%%%%%%%%%%%%%%%%%%%%%%%%%%%
\appendix
%%%%%%%%%%%%%%%%%%%%%%%%%%%%%%%%%%%%%%%%%%%%%%%%%%%%%%%%%%%%%%%%%%%%%%%%%%%

\section{Experimental Parameters}
\label{app:params}

For reproducibility, we record the exact parameters of all reported experiments.

\begin{table}[htbp]
\centering
\caption{Parameters common to all experiments.}
\renewcommand{\arraystretch}{1.2}
\begin{tabular}{@{}ll@{}}
\toprule
Parameter & Value \\
\midrule
Viscosity $\nu$ & $10^{-3}$ \\
Domain & $[0, 2\pi]^d$ (periodic) \\
Dealiasing & 2/3 rule (all runs) \\
Time integration & RK4 (momentum, outer loop) \\
$\theta$ time integration & IMEX: CN implicit, RK4 explicit \\
CFL condition & $\mathrm{CFL} = 0.5$ (adaptive) \\
Random seed & 42 (all runs) \\
Results database & \texttt{results/results.db} (SQLite) \\
\midrule
2D grid size & $N = 64^2$ \\
2D run time & $T = 10$ \\
3D grid sizes & $N \in \{32^3, 64^3, 128^3, 256^3\}$ \\
3D hardware & Apple M5 (10-core), JAX 0.4.26, jax-metal 0.1.0 \\
\bottomrule
\end{tabular}
\end{table}

\section{The Wang et al.\ Self-Similar Profiles}
\label{app:wang}

The self-similar profiles used in Sections~\ref{subsec:3D} and~\ref{subsec:blowup}
are parameterised by $\lambda \in \R$ following Wang~\cite{Wang2025}. The
exact values used in our experiments are:

\begin{center}
\renewcommand{\arraystretch}{1.2}
\begin{tabular}{@{}lll@{}}
\toprule
Profile & $\lambda$ & Type \\
\midrule
CCF stable & $1.1808$ & Stable \\
CCF 1st unstable & $0.6057$ & Unstable (critical) \\
CCF 2nd unstable & $0.4703$ & Unstable (most dangerous) \\
Boussinesq stable & $1.9206$ & Stable \\
Boussinesq 1st & $1.3991$ & Unstable \\
Boussinesq 3rd & $1.1843$ & Unstable \\
Adversarial minimum & $0.05$ & Extremal (near $\lambda = 0$) \\
\bottomrule
\end{tabular}
\end{center}

The 3D embedding proceeds by constructing vortex tubes aligned with the $z$-axis
using the 2D Wang et al.\ vorticity profile $\omega_{2D}(r, \lambda)$ and
adding a small $z$-perturbation to break two-dimensionality.
The embedding is implemented in \texttt{layer3/wang\_profiles\_3D.py}.

\end{document}
